% Gerado a partir de docs/tcc.md por scripts/sincronizar_tcc.py.
\chapter{Introdução}\label{sec:1}

\section{Contextualização}\label{sec:1-1}

Nas últimas décadas, a relação entre Estado e cidadão passou por uma mudança estrutural: atividades que antes dependiam de balcões presenciais, deslocamentos e atendimento físico foram progressivamente incorporadas a portais, aplicações e fluxos digitais. A plataforma GOV.BR simboliza esse movimento ao concentrar a oferta digital de serviços federais e ao reforçar a ideia de que governo digital não é apenas troca de canal, mas reorganização da forma como o cidadão encontra, solicita, acompanha e avalia serviços públicos.

Os serviços públicos compõem uma das formas mais diretas dessa relação. Solicitar um documento, consultar uma situação cadastral, acompanhar um benefício, registrar uma demanda ou acessar uma política pública são situações em que a experiência do usuário influencia a qualidade percebida da administração pública. A digitalização ampliou o alcance e a velocidade desses atendimentos, mas também deslocou parte das dificuldades para a própria navegação: localizar o serviço correto, entender a linguagem institucional, preencher etapas intermediárias e perceber se a solicitação realmente avançou.

A escala dessa transição aparece de forma clara em experiências recentes de participação digital. O Plano Plurianual (PPA) Participativo 2024-2027, por exemplo, mobilizou mais de 1,5 milhão de participantes e superou 4 milhões de acessos na plataforma Brasil Participativo, demonstrando a capacidade de mobilização do ambiente digital. Ao mesmo tempo, a experiência registrou barreiras como problemas de conectividade, dificuldade de uso, fluxos de cadastro complexos e necessidade de melhor retorno ao cidadão sobre sua contribuição \cite{brasilMgi2024ppaParticipativo}. Esse contraste reforça a importância de avaliar não apenas se o serviço existe em formato digital, mas como ele é percebido e utilizado durante a jornada.

No Brasil, a avaliação dos serviços públicos possui base legal e gerencial. A Lei n\textsuperscript{o} 13.460/2017 trata da participação, proteção e defesa dos usuários de serviços públicos e prevê que os resultados de avaliação subsidiem a reorientação e o ajuste da prestação estatal \cite{brasil2017lei13460}. O Decreto n\textsuperscript{o} 9.094/2017 institui a Carta de Serviços ao Usuário e reforça a simplificação do atendimento \cite{brasil2017decreto9094}. A Lei n\textsuperscript{o} 14.129/2021 associa governo digital, eficiência pública, desburocratização e participação do cidadão \cite{brasil2021lei14129}. Como referência histórica, o Decreto n\textsuperscript{o} 11.260/2022 tratou da elaboração da Estratégia Nacional de Governo Digital e da prorrogação da estratégia anterior; foi revogado em 2024 e não é apresentado como norma vigente \cite{brasil2022decreto11260}.

A avaliação da experiência digital, porém, não depende apenas de perguntar ao usuário se ele ficou satisfeito ao final do atendimento. Em serviços digitais, excesso de cliques, permanência por tempo acima do esperado, retorno repetido a páginas anteriores e abandono antes da conclusão constituem evidências comportamentais de possível dificuldade. Esses sinais não medem qualidade percebida e não autorizam, isoladamente, concluir que a experiência foi ruim. Quando coletados com finalidade clara e com cuidado em relação à privacidade, eles ajudam a localizar pontos da jornada que devem ser investigados e confrontados com a percepção declarada pelo usuário \cite{deloneMclean2003,govuk2021completionRate,brasil2018lgpd}.

A proposta é relacionar evidências comportamentais e declaradas sem tratá-las como equivalentes. A implementação Python registra cliques em elementos instrumentados, página, serviço e informações do evento; recebe também uma duração informada à API, demonstrada por envio simulado de 120 segundos, sem cronômetro automático no cliente. Regras operacionais podem gerar uma sugestão de redirecionamento ou disponibilizar um convite. Estrelas, atributos, CES e comentários permitem examinar julgamentos declarados sobre satisfação, aspectos do serviço, facilidade percebida e contexto qualitativo, respectivamente. Esses indicadores têm interpretações distintas e não formam um escore único de qualidade. O estudo com participantes ocorreu separadamente, em interfaces no Figma, com cliques e tempos registrados automaticamente pelo Useberry.

A proposta adota uma interface curta de avaliação e mantém uma implementação local de registro de eventos e feedback. A visualização administrativa disponível apresenta alertas, contagens por serviço e médias de respostas. O painel apresenta frequência de atributos, sessões distintas alertadas e ações com acompanhamento. Taxa de resposta por fluxo ainda depende de instrumentação A/B e de exposição. A Lean Inception orientou a delimitação do produto e do MVP, cujas funcionalidades foram implementadas parcialmente \cite{caroli2018leanInception,fowler2022leanInception}.

\section{Problema de Pesquisa}\label{sec:1-2}

A questão investigada é como relacionar a avaliação declarada pelo cidadão ao contexto de navegação em que ela é solicitada. Quando a coleta fica restrita ao final da jornada, pode deixar de registrar detalhes de uma etapa anterior ou de um fluxo interrompido. Essa possibilidade motiva o estudo do momento do convite, sem pressupor que a coleta contextual produza resultados superiores. Indicadores só se tornam úteis à gestão quando seu contexto e suas limitações são considerados \cite{govuk2021usingData,deloneMclean2003}.

Modelos oficiais e instrumentos consolidados de avaliação de serviços públicos contribuem para a mensuração da satisfação e da qualidade percebida. Ainda assim, permanecem desafios relacionados à baixa adesão, ao esforço exigido do usuário e à transformação das respostas em informação acionável para o gestor. Avaliações longas podem reduzir a participação; avaliações curtas podem gerar indicadores pouco explicativos; dados de navegação sem interpretação podem produzir apenas registros técnicos sem valor gerencial. A literatura científica mostra que taxas de resposta variam entre modos de coleta e populações, e que a formulação e o contexto das perguntas influenciam os relatos dos respondentes \cite{shihFan2008responseRates,manfreda2008webSurveys,schwarz1999selfReports}. Esses estudos justificam examinar adesão e desenho da coleta, mas não demonstram superioridade do feedback contextual em serviços públicos digitais. O uso dos registros para melhoria também exige interpretação dos dados de desempenho e atenção à relação entre qualidade do sistema, satisfação e benefícios \cite{govuk2021usingData,deloneMclean2003}.

Este trabalho parte da seguinte questão de pesquisa: como propor e prototipar um avaliador de serviços públicos digitais que identifique sinais de possível dificuldade na navegação, colete feedback contextual de baixo esforço e organize informações úteis para a melhoria dos serviços avaliados?

A comparação entre convite contextual e formulário final foi organizada nas hipóteses H1 e H2, formuladas na seção \ref{sec:3-6-1}. Elas tratam, respectivamente, de participação e especificidade dos comentários do formulário. Os resultados disponíveis são observações exploratórias de interface e não testam essas hipóteses.

\section{Objetivo Geral}\label{sec:1-3}

Propor e desenvolver um protótipo inicial de avaliador de serviços públicos digitais capaz de registrar sinais de possível dificuldade na navegação, coletar feedback contextual estruturado, comparar esse fluxo com uma avaliação estática e apresentar indicadores administrativos que apoiem a melhoria contínua dos serviços.

\section{Objetivos Específicos}\label{sec:1-4}

Para alcançar o objetivo geral, foram definidos os seguintes objetivos específicos:

\begin{itemize}
\item
  levantar fundamentos legais, técnicos e teóricos sobre avaliação de serviços públicos, qualidade percebida, usabilidade, governo digital e Lean Inception;
\item
  definir, por meio de Lean Inception, o fluxo de detecção de possível dificuldade, sugestão de redirecionamento e coleta de feedback;
\item
  construir um protótipo funcional com API, páginas simuladas de serviço público, feedback contextual, formulário estático comparativo e visualização administrativa;
\item
  estruturar os dados e indicadores necessários para apoiar a análise de cliques, tempo de navegação, taxa de resposta, satisfação por estrelas, esforço percebido pelo CES e pontos recorrentes de melhoria.
\end{itemize}

A estrutura dos objetivos combina a fundamentação legal e técnica da avaliação de serviços públicos com práticas de definição de MVP e prototipação de software \cite{brasil2017lei13460,brasil2021lei14129,caroli2018leanInception,Sommerville2011}.

\section{Justificativa}\label{sec:1-5}

Mecanismos digitais de avaliação podem fortalecer a participação do usuário na melhoria dos serviços públicos. Os materiais da Enap destacam que a avaliação centrada no cidadão desloca o foco das ferramentas tecnológicas para as necessidades, demandas e satisfação dos usuários \cite{limaPereira2020modulo1}. Também indicam que avaliações devem alimentar ciclos de melhoria, permitindo identificar problemas recorrentes, atributos críticos e prioridades de intervenção \cite{limaPereira2020modulo4}.

A proposta é relevante porque muitos usuários já conhecem avaliações simples por estrelas, motivos e comentários. Uma solução pública pode aproveitar essa familiaridade para reduzir o esforço de resposta, desde que use os dados para finalidades compatíveis com o interesse público. Além disso, a combinação entre eventos de navegação e feedback declarado permite que o gestor observe tanto o que o usuário informou quanto o que ocorreu durante a jornada digital.

Do ponto de vista investigativo, o trabalho formula hipóteses sobre o momento da coleta e distingue indícios comportamentais de percepção declarada. A organização conceitual é explicitada no Quadro \ref{tab:construtos-indicadores-dados}, com dados, origem e limites de interpretação; sua validade ainda depende de avaliação empírica. Do ponto de vista tecnológico, a contribuição é a implementação de API, telas, regras explícitas e painel que permitem associar contexto de navegação e relato do cidadão. As configurações possuem versões e histórico local, o que permite inspecionar alterações, sem equivaler a auditoria independente \cite{deloneMclean2003,hevner2004designscience,Sommerville2011}.

\section{Escopo Inicial}\label{sec:1-6}

O escopo inicial concentra-se na proposta e prototipação de uma solução. O trabalho não pretende implantar um sistema em ambiente governamental real, nem substituir instrumentos oficiais de avaliação já existentes. A solução é complementar e exploratória: simula páginas de serviço público, registra eventos de navegação, aplica regras configuráveis de acionamento por sinais de possível dificuldade, abre uma coleta simples de feedback, inclui um formulário reutilizável e apresenta dados administrativos agregados. A comparação entre formulário final e convite contextual pertence ao desenho de interface no Figma; a aplicação Python ainda não distingue os grupos nem registra a exposição. Essa delimitação segue a lógica de prototipação como forma de reduzir incertezas antes de uma solução definitiva e mantém cuidado com minimização de dados pessoais \cite{Sommerville2011,brasil2018lgpd}.

A avaliação desta etapa reúne procedimentos distintos: inspeção funcional e heurística interna/preliminar relatada pelos autores, e sessões exploratórias de interface com oito participantes no Figma, instrumentadas pelo Useberry. Essas sessões não constituem teste de integração ou execução da API Python. Os resultados são descritivos, sem generalização estatística ou validação em contexto real de prestação de serviço público. Integração com portais governamentais, calibração e estudo comparativo permanecem etapas futuras.
